% =========================================================================
% Bibliografia del TFG -- Pablo Chantada Saborido
%
% Uso:
%   (a) Compilar este fichero por si solo para previsualizar:
%       pdflatex bibliografia.tex
%   (b) Incluirlo desde la memoria principal:
%       % =========================================================================
% Bibliografia del TFG -- Pablo Chantada Saborido
%
% Uso:
%   (a) Compilar este fichero por si solo para previsualizar:
%       pdflatex bibliografia.tex
%   (b) Incluirlo desde la memoria principal:
%       % =========================================================================
% Bibliografia del TFG -- Pablo Chantada Saborido
%
% Uso:
%   (a) Compilar este fichero por si solo para previsualizar:
%       pdflatex bibliografia.tex
%   (b) Incluirlo desde la memoria principal:
%       % =========================================================================
% Bibliografia del TFG -- Pablo Chantada Saborido
%
% Uso:
%   (a) Compilar este fichero por si solo para previsualizar:
%       pdflatex bibliografia.tex
%   (b) Incluirlo desde la memoria principal:
%       \input{bibliografia}   % dentro del documento, antes de \end{document}
%   (c) Copiar solo el bloque \begin{thebibliography}...\end{thebibliography}
%       al final de tu memoria.tex
%
% Estilo: numerico tipo IEEE. Los \cite{key} usan las claves indicadas.
% Entradas marcadas con (*) son las imprescindibles.
% =========================================================================

\documentclass[11pt,a4paper]{article}
\usepackage[utf8]{inputenc}
\usepackage[T1]{fontenc}
\usepackage[spanish]{babel}
\usepackage{geometry}
\usepackage{hyperref}
\geometry{margin=2.5cm}

\title{Bibliografia --- TFG Agentes RL en Minecraft}
\author{Pablo Chantada Saborido}
\date{}

\begin{document}
\maketitle

\begin{thebibliography}{99}

% -------------------------------------------------------------------------
% 1. Plataformas y entornos Minecraft
% -------------------------------------------------------------------------

\bibitem{malmo2016}
M.~Johnson, K.~Hofmann, T.~Hutton, and D.~Bignell,
``The Malmo Platform for Artificial Intelligence Experimentation,''
in \emph{Proc.\ 25th Int.\ Joint Conf.\ on Artificial Intelligence (IJCAI)},
2016, pp.~4246--4247.

\bibitem{minerl2019}
W.~H. Guss, B.~Houghton, N.~Topin, P.~Wang, C.~Codel, M.~Veloso, and
R.~Salakhutdinov,
``MineRL: A Large-Scale Dataset of Minecraft Demonstrations,''
in \emph{Proc.\ 28th Int.\ Joint Conf.\ on Artificial Intelligence (IJCAI)},
2019, pp.~2442--2448.

\bibitem{minerl2020}
W.~H. Guss \emph{et al.},
``Towards Robust and Domain Agnostic Reinforcement Learning Competitions:
MineRL 2020,''
in \emph{Proc.\ NeurIPS 2020 Competition and Demonstration Track},
PMLR, vol.~133, 2021, pp.~233--252.

\bibitem{minerldiamond2021}
A.~Kanervisto \emph{et al.},
``MineRL Diamond 2021 Competition: Overview, Results, and Lessons Learned,''
in \emph{NeurIPS 2021 Competitions and Demonstrations Track},
2022.

\bibitem{minedojo2022}
L.~Fan \emph{et al.},
``MineDojo: Building Open-Ended Embodied Agents with Internet-Scale
Knowledge,''
in \emph{Advances in Neural Information Processing Systems (NeurIPS)},
vol.~35, 2022.

\bibitem{mineflayer}
PrismarineJS,
\emph{Mineflayer: Create Minecraft bots with a high-level JavaScript API}
[Software], 2024. [Online]. Available:
\url{https://github.com/PrismarineJS/mineflayer}

\bibitem{papermc}
PaperMC,
\emph{Paper: High performance Minecraft server} [Software], 2024.
[Online]. Available: \url{https://papermc.io/}

% -------------------------------------------------------------------------
% 2. Agentes con LLMs en Minecraft
% -------------------------------------------------------------------------

\bibitem{voyager2023}
G.~Wang, Y.~Xie, Y.~Jiang, A.~Mandlekar, C.~Xiao, Y.~Zhu, L.~Fan, and
A.~Anandkumar,
``Voyager: An Open-Ended Embodied Agent with Large Language Models,''
\emph{arXiv preprint arXiv:2305.16291}, 2023.

\bibitem{ghost2023}
X.~Zhu \emph{et al.},
``Ghost in the Minecraft: Generally Capable Agents for Open-World Environments
via Large Language Models with Text-based Knowledge and Memory,''
\emph{arXiv preprint arXiv:2305.17144}, 2023.

\bibitem{jarvis1}
Z.~Wang \emph{et al.},
``JARVIS-1: Open-World Multi-task Agents with Memory-Augmented Multimodal
Language Models,''
\emph{arXiv preprint arXiv:2311.05997}, 2023.

\bibitem{steve1}
S.~Lifshitz, K.~Paster, H.~Chan, J.~Ba, and S.~McIlraith,
``STEVE-1: A Generative Model for Text-to-Behavior in Minecraft,''
in \emph{Advances in Neural Information Processing Systems (NeurIPS)},
2023.

\bibitem{plan4mc}
H.~Yuan, C.~Zhang, H.~Wang, F.~Xie, P.~Cai, H.~Dong, and Z.~Lu,
``Plan4MC: Skill Reinforcement Learning and Planning for Open-World Minecraft
Tasks,''
\emph{arXiv preprint arXiv:2303.16563}, 2023.

\bibitem{mindcraft}
K.~Nottingham \emph{et al.},
\emph{Mindcraft: LLM-powered Minecraft agents} [Software], 2024.
[Online]. Available: \url{https://github.com/kolbytn/mindcraft}

% -------------------------------------------------------------------------
% 3. Planificacion clasica (HTN)
% -------------------------------------------------------------------------

\bibitem{erol1994}
K.~Erol, J.~Hendler, and D.~S. Nau,
``HTN Planning: Complexity and Expressivity,''
in \emph{Proc.\ 12th National Conf.\ on Artificial Intelligence (AAAI)},
1994, pp.~1123--1128.

\bibitem{ghallab2004}
M.~Ghallab, D.~Nau, and P.~Traverso,
\emph{Automated Planning: Theory and Practice}.
San Francisco, CA: Morgan Kaufmann, 2004.

\bibitem{shop2}
D.~S. Nau, T.-C. Au, O.~Ilghami, U.~Kuter, J.~W. Murdock, D.~Wu, and F.~Yaman,
``SHOP2: An HTN Planning System,''
\emph{Journal of Artificial Intelligence Research}, vol.~20, pp.~379--404,
2003.

\bibitem{georgievski2015}
I.~Georgievski and M.~Aiello,
``HTN planning: Overview, comparison, and beyond,''
\emph{Artificial Intelligence}, vol.~222, pp.~124--156, 2015.

% -------------------------------------------------------------------------
% 4. Imitation Learning y Behavior Cloning
% -------------------------------------------------------------------------

\bibitem{pomerleau1991}
D.~A. Pomerleau,
``Efficient Training of Artificial Neural Networks for Autonomous
Navigation,''
\emph{Neural Computation}, vol.~3, no.~1, pp.~88--97, 1991.

\bibitem{dagger}
S.~Ross, G.~Gordon, and D.~Bagnell,
``A Reduction of Imitation Learning and Structured Prediction to No-Regret
Online Learning,''
in \emph{Proc.\ 14th Int.\ Conf.\ on Artificial Intelligence and Statistics
(AISTATS)}, 2011, pp.~627--635.

\bibitem{gail}
J.~Ho and S.~Ermon,
``Generative Adversarial Imitation Learning,''
in \emph{Advances in Neural Information Processing Systems (NeurIPS)},
vol.~29, 2016.

\bibitem{hussein2017}
A.~Hussein, M.~M. Gaber, E.~Elyan, and C.~Jayne,
``Imitation Learning: A Survey of Learning Methods,''
\emph{ACM Computing Surveys}, vol.~50, no.~2, art.~21, 2017.

\bibitem{vpt2022}
B.~Baker \emph{et al.},
``Video PreTraining (VPT): Learning to Act by Watching Unlabeled Online
Videos,''
in \emph{Advances in Neural Information Processing Systems (NeurIPS)},
vol.~35, 2022.

% -------------------------------------------------------------------------
% 5. Deep Reinforcement Learning -- fundamentos
% -------------------------------------------------------------------------

\bibitem{suttonbarto}
R.~S. Sutton and A.~G. Barto,
\emph{Reinforcement Learning: An Introduction}, 2nd ed.
Cambridge, MA: MIT Press, 2018.

\bibitem{dqnnature}
V.~Mnih \emph{et al.},
``Human-level control through deep reinforcement learning,''
\emph{Nature}, vol.~518, pp.~529--533, 2015.

\bibitem{atari2013}
V.~Mnih, K.~Kavukcuoglu, D.~Silver, A.~Graves, I.~Antonoglou, D.~Wierstra,
and M.~Riedmiller,
``Playing Atari with Deep Reinforcement Learning,''
\emph{arXiv preprint arXiv:1312.5602}, 2013.

\bibitem{doubledqn}
H.~van Hasselt, A.~Guez, and D.~Silver,
``Deep Reinforcement Learning with Double Q-Learning,''
in \emph{Proc.\ 30th AAAI Conf.\ on Artificial Intelligence}, 2016,
pp.~2094--2100.

\bibitem{dueling}
Z.~Wang, T.~Schaul, M.~Hessel, H.~van Hasselt, M.~Lanctot, and N.~de~Freitas,
``Dueling Network Architectures for Deep Reinforcement Learning,''
in \emph{Proc.\ 33rd Int.\ Conf.\ on Machine Learning (ICML)}, 2016.

\bibitem{per}
T.~Schaul, J.~Quan, I.~Antonoglou, and D.~Silver,
``Prioritized Experience Replay,''
in \emph{Proc.\ Int.\ Conf.\ on Learning Representations (ICLR)}, 2016.

\bibitem{rainbow}
M.~Hessel \emph{et al.},
``Rainbow: Combining Improvements in Deep Reinforcement Learning,''
in \emph{Proc.\ 32nd AAAI Conf.\ on Artificial Intelligence}, 2018.

% -------------------------------------------------------------------------
% 6. Policy gradient: PPO, SAC y variantes
% -------------------------------------------------------------------------

\bibitem{ppo}
J.~Schulman, F.~Wolski, P.~Dhariwal, A.~Radford, and O.~Klimov,
``Proximal Policy Optimization Algorithms,''
\emph{arXiv preprint arXiv:1707.06347}, 2017.

\bibitem{sac2018}
T.~Haarnoja, A.~Zhou, P.~Abbeel, and S.~Levine,
``Soft Actor-Critic: Off-Policy Maximum Entropy Deep Reinforcement Learning
with a Stochastic Actor,''
in \emph{Proc.\ 35th Int.\ Conf.\ on Machine Learning (ICML)}, 2018.

\bibitem{sacapp}
T.~Haarnoja \emph{et al.},
``Soft Actor-Critic Algorithms and Applications,''
\emph{arXiv preprint arXiv:1812.05905}, 2018.

\bibitem{sacdiscrete}
P.~Christodoulou,
``Soft Actor-Critic for Discrete Action Settings,''
\emph{arXiv preprint arXiv:1910.07207}, 2019.

\bibitem{a3c}
V.~Mnih \emph{et al.},
``Asynchronous Methods for Deep Reinforcement Learning,''
in \emph{Proc.\ 33rd Int.\ Conf.\ on Machine Learning (ICML)}, 2016.

\bibitem{andrychowicz2021}
M.~Andrychowicz \emph{et al.},
``What Matters for On-Policy Deep Actor-Critic Methods? A Large-Scale
Study,''
in \emph{Proc.\ Int.\ Conf.\ on Learning Representations (ICLR)}, 2021.

% -------------------------------------------------------------------------
% 7. RL + demostraciones (Hybrid, BC + RL)
% -------------------------------------------------------------------------

\bibitem{dqfd}
T.~Hester \emph{et al.},
``Deep Q-Learning from Demonstrations,''
in \emph{Proc.\ 32nd AAAI Conf.\ on Artificial Intelligence}, 2018.

\bibitem{ddpgfd}
M.~Vecerik \emph{et al.},
``Leveraging Demonstrations for Deep Reinforcement Learning on Robotics
Problems with Sparse Rewards,''
\emph{arXiv preprint arXiv:1707.08817}, 2017.

\bibitem{nair2018}
A.~Nair, B.~McGrew, M.~Andrychowicz, W.~Zaremba, and P.~Abbeel,
``Overcoming Exploration in Reinforcement Learning with Demonstrations,''
in \emph{Proc.\ IEEE Int.\ Conf.\ on Robotics and Automation (ICRA)}, 2018.

\bibitem{dapg}
A.~Rajeswaran \emph{et al.},
``Learning Complex Dexterous Manipulation with Deep Reinforcement Learning
and Demonstrations,''
in \emph{Robotics: Science and Systems (RSS)}, 2018.

\bibitem{sqil}
S.~Reddy, A.~D. Dragan, and S.~Levine,
``SQIL: Imitation Learning via Reinforcement Learning with Sparse Rewards,''
in \emph{Proc.\ Int.\ Conf.\ on Learning Representations (ICLR)}, 2020.

\bibitem{iqlearn}
D.~Garg, S.~Chakraborty, C.~Cundy, J.~Song, and S.~Ermon,
``IQ-Learn: Inverse soft-Q Learning for Imitation,''
in \emph{Advances in Neural Information Processing Systems (NeurIPS)},
vol.~34, 2021.

% -------------------------------------------------------------------------
% 8. RL visual, CNN y normalizacion
% -------------------------------------------------------------------------

\bibitem{resnet}
K.~He, X.~Zhang, S.~Ren, and J.~Sun,
``Deep Residual Learning for Image Recognition,''
in \emph{Proc.\ IEEE Conf.\ on Computer Vision and Pattern Recognition
(CVPR)}, 2016, pp.~770--778.

\bibitem{groupnorm}
Y.~Wu and K.~He,
``Group Normalization,''
in \emph{Proc.\ European Conf.\ on Computer Vision (ECCV)}, 2018,
pp.~3--19.

\bibitem{impala}
L.~Espeholt \emph{et al.},
``IMPALA: Scalable Distributed Deep-RL with Importance Weighted
Actor-Learner Architectures,''
in \emph{Proc.\ 35th Int.\ Conf.\ on Machine Learning (ICML)}, 2018.

\bibitem{drqv2}
D.~Yarats, R.~Fergus, A.~Lazaric, and L.~Pinto,
``Mastering Visual Continuous Control: Improved Data-Augmented
Reinforcement Learning,''
\emph{arXiv preprint arXiv:2107.09645}, 2021.

% -------------------------------------------------------------------------
% 9. World models y RL jerarquico (lineas futuras)
% -------------------------------------------------------------------------

\bibitem{worldmodels}
D.~Ha and J.~Schmidhuber,
``World Models,''
in \emph{Advances in Neural Information Processing Systems (NeurIPS)},
2018.

\bibitem{dreamerv3}
D.~Hafner, J.~Pasukonis, J.~Ba, and T.~Lillicrap,
``Mastering Diverse Domains through World Models,''
\emph{arXiv preprint arXiv:2301.04104}, 2023.

\bibitem{muzero}
J.~Schrittwieser \emph{et al.},
``Mastering Atari, Go, Chess and Shogi by Planning with a Learned Model,''
\emph{Nature}, vol.~588, pp.~604--609, 2020.

\bibitem{options}
R.~S. Sutton, D.~Precup, and S.~Singh,
``Between MDPs and semi-MDPs: A Framework for Temporal Abstraction in
Reinforcement Learning,''
\emph{Artificial Intelligence}, vol.~112, no.~1--2, pp.~181--211, 1999.

% -------------------------------------------------------------------------
% 10. Multi-agente
% -------------------------------------------------------------------------

\bibitem{coma}
J.~Foerster, G.~Farquhar, T.~Afouras, N.~Nardelli, and S.~Whiteson,
``Counterfactual Multi-Agent Policy Gradients,''
in \emph{Proc.\ 32nd AAAI Conf.\ on Artificial Intelligence}, 2018.

\bibitem{qmix}
T.~Rashid, M.~Samvelyan, C.~Schroeder de~Witt, G.~Farquhar, J.~Foerster,
and S.~Whiteson,
``QMIX: Monotonic Value Function Factorisation for Deep Multi-Agent
Reinforcement Learning,''
in \emph{Proc.\ 35th Int.\ Conf.\ on Machine Learning (ICML)}, 2018.

\bibitem{maddpg}
R.~Lowe, Y.~Wu, A.~Tamar, J.~Harb, P.~Abbeel, and I.~Mordatch,
``Multi-Agent Actor-Critic for Mixed Cooperative-Competitive
Environments,''
in \emph{Advances in Neural Information Processing Systems (NeurIPS)},
vol.~30, 2017.

% -------------------------------------------------------------------------
% 11. Vision por computador (YOLO, ResNet -- si justificas pipeline visual)
% -------------------------------------------------------------------------

\bibitem{yolo2016}
J.~Redmon, S.~Divvala, R.~Girshick, and A.~Farhadi,
``You Only Look Once: Unified, Real-Time Object Detection,''
in \emph{Proc.\ IEEE Conf.\ on Computer Vision and Pattern Recognition
(CVPR)}, 2016, pp.~779--788.

\bibitem{ultralytics}
G.~Jocher \emph{et al.},
\emph{Ultralytics YOLO} [Software], 2024.
[Online]. Available: \url{https://github.com/ultralytics/ultralytics}

\bibitem{alexnet}
A.~Krizhevsky, I.~Sutskever, and G.~E. Hinton,
``ImageNet Classification with Deep Convolutional Neural Networks,''
in \emph{Advances in Neural Information Processing Systems (NeurIPS)},
vol.~25, 2012, pp.~1097--1105.

\bibitem{convlstm}
X.~Shi, Z.~Chen, H.~Wang, D.-Y. Yeung, W.-K. Wong, and W.-C. Woo,
``Convolutional LSTM Network: A Machine Learning Approach for Precipitation
Nowcasting,''
in \emph{Advances in Neural Information Processing Systems (NeurIPS)},
vol.~28, 2015.

% -------------------------------------------------------------------------
% 12. Surveys y panoramicas (estado del arte)
% -------------------------------------------------------------------------

\bibitem{arulkumaran2017}
K.~Arulkumaran, M.~P. Deisenroth, M.~Brundage, and A.~A. Bharath,
``Deep Reinforcement Learning: A Brief Survey,''
\emph{IEEE Signal Processing Magazine}, vol.~34, no.~6, pp.~26--38, 2017.

\bibitem{li2018}
Y.~Li,
``Deep Reinforcement Learning: An Overview,''
\emph{arXiv preprint arXiv:1701.07274}, 2018.

\bibitem{agentdivisions}
B.~Liu \emph{et al.},
``A Survey of Agent Divisions and Cooperative Frameworks in Multi-Agent
Reinforcement Learning,''
\emph{arXiv preprint arXiv:2111.08857}, 2021.

\end{thebibliography}

\end{document}
   % dentro del documento, antes de \end{document}
%   (c) Copiar solo el bloque \begin{thebibliography}...\end{thebibliography}
%       al final de tu memoria.tex
%
% Estilo: numerico tipo IEEE. Los \cite{key} usan las claves indicadas.
% Entradas marcadas con (*) son las imprescindibles.
% =========================================================================

\documentclass[11pt,a4paper]{article}
\usepackage[utf8]{inputenc}
\usepackage[T1]{fontenc}
\usepackage[spanish]{babel}
\usepackage{geometry}
\usepackage{hyperref}
\geometry{margin=2.5cm}

\title{Bibliografia --- TFG Agentes RL en Minecraft}
\author{Pablo Chantada Saborido}
\date{}

\begin{document}
\maketitle

\begin{thebibliography}{99}

% -------------------------------------------------------------------------
% 1. Plataformas y entornos Minecraft
% -------------------------------------------------------------------------

\bibitem{malmo2016}
M.~Johnson, K.~Hofmann, T.~Hutton, and D.~Bignell,
``The Malmo Platform for Artificial Intelligence Experimentation,''
in \emph{Proc.\ 25th Int.\ Joint Conf.\ on Artificial Intelligence (IJCAI)},
2016, pp.~4246--4247.

\bibitem{minerl2019}
W.~H. Guss, B.~Houghton, N.~Topin, P.~Wang, C.~Codel, M.~Veloso, and
R.~Salakhutdinov,
``MineRL: A Large-Scale Dataset of Minecraft Demonstrations,''
in \emph{Proc.\ 28th Int.\ Joint Conf.\ on Artificial Intelligence (IJCAI)},
2019, pp.~2442--2448.

\bibitem{minerl2020}
W.~H. Guss \emph{et al.},
``Towards Robust and Domain Agnostic Reinforcement Learning Competitions:
MineRL 2020,''
in \emph{Proc.\ NeurIPS 2020 Competition and Demonstration Track},
PMLR, vol.~133, 2021, pp.~233--252.

\bibitem{minerldiamond2021}
A.~Kanervisto \emph{et al.},
``MineRL Diamond 2021 Competition: Overview, Results, and Lessons Learned,''
in \emph{NeurIPS 2021 Competitions and Demonstrations Track},
2022.

\bibitem{minedojo2022}
L.~Fan \emph{et al.},
``MineDojo: Building Open-Ended Embodied Agents with Internet-Scale
Knowledge,''
in \emph{Advances in Neural Information Processing Systems (NeurIPS)},
vol.~35, 2022.

\bibitem{mineflayer}
PrismarineJS,
\emph{Mineflayer: Create Minecraft bots with a high-level JavaScript API}
[Software], 2024. [Online]. Available:
\url{https://github.com/PrismarineJS/mineflayer}

\bibitem{papermc}
PaperMC,
\emph{Paper: High performance Minecraft server} [Software], 2024.
[Online]. Available: \url{https://papermc.io/}

% -------------------------------------------------------------------------
% 2. Agentes con LLMs en Minecraft
% -------------------------------------------------------------------------

\bibitem{voyager2023}
G.~Wang, Y.~Xie, Y.~Jiang, A.~Mandlekar, C.~Xiao, Y.~Zhu, L.~Fan, and
A.~Anandkumar,
``Voyager: An Open-Ended Embodied Agent with Large Language Models,''
\emph{arXiv preprint arXiv:2305.16291}, 2023.

\bibitem{ghost2023}
X.~Zhu \emph{et al.},
``Ghost in the Minecraft: Generally Capable Agents for Open-World Environments
via Large Language Models with Text-based Knowledge and Memory,''
\emph{arXiv preprint arXiv:2305.17144}, 2023.

\bibitem{jarvis1}
Z.~Wang \emph{et al.},
``JARVIS-1: Open-World Multi-task Agents with Memory-Augmented Multimodal
Language Models,''
\emph{arXiv preprint arXiv:2311.05997}, 2023.

\bibitem{steve1}
S.~Lifshitz, K.~Paster, H.~Chan, J.~Ba, and S.~McIlraith,
``STEVE-1: A Generative Model for Text-to-Behavior in Minecraft,''
in \emph{Advances in Neural Information Processing Systems (NeurIPS)},
2023.

\bibitem{plan4mc}
H.~Yuan, C.~Zhang, H.~Wang, F.~Xie, P.~Cai, H.~Dong, and Z.~Lu,
``Plan4MC: Skill Reinforcement Learning and Planning for Open-World Minecraft
Tasks,''
\emph{arXiv preprint arXiv:2303.16563}, 2023.

\bibitem{mindcraft}
K.~Nottingham \emph{et al.},
\emph{Mindcraft: LLM-powered Minecraft agents} [Software], 2024.
[Online]. Available: \url{https://github.com/kolbytn/mindcraft}

% -------------------------------------------------------------------------
% 3. Planificacion clasica (HTN)
% -------------------------------------------------------------------------

\bibitem{erol1994}
K.~Erol, J.~Hendler, and D.~S. Nau,
``HTN Planning: Complexity and Expressivity,''
in \emph{Proc.\ 12th National Conf.\ on Artificial Intelligence (AAAI)},
1994, pp.~1123--1128.

\bibitem{ghallab2004}
M.~Ghallab, D.~Nau, and P.~Traverso,
\emph{Automated Planning: Theory and Practice}.
San Francisco, CA: Morgan Kaufmann, 2004.

\bibitem{shop2}
D.~S. Nau, T.-C. Au, O.~Ilghami, U.~Kuter, J.~W. Murdock, D.~Wu, and F.~Yaman,
``SHOP2: An HTN Planning System,''
\emph{Journal of Artificial Intelligence Research}, vol.~20, pp.~379--404,
2003.

\bibitem{georgievski2015}
I.~Georgievski and M.~Aiello,
``HTN planning: Overview, comparison, and beyond,''
\emph{Artificial Intelligence}, vol.~222, pp.~124--156, 2015.

% -------------------------------------------------------------------------
% 4. Imitation Learning y Behavior Cloning
% -------------------------------------------------------------------------

\bibitem{pomerleau1991}
D.~A. Pomerleau,
``Efficient Training of Artificial Neural Networks for Autonomous
Navigation,''
\emph{Neural Computation}, vol.~3, no.~1, pp.~88--97, 1991.

\bibitem{dagger}
S.~Ross, G.~Gordon, and D.~Bagnell,
``A Reduction of Imitation Learning and Structured Prediction to No-Regret
Online Learning,''
in \emph{Proc.\ 14th Int.\ Conf.\ on Artificial Intelligence and Statistics
(AISTATS)}, 2011, pp.~627--635.

\bibitem{gail}
J.~Ho and S.~Ermon,
``Generative Adversarial Imitation Learning,''
in \emph{Advances in Neural Information Processing Systems (NeurIPS)},
vol.~29, 2016.

\bibitem{hussein2017}
A.~Hussein, M.~M. Gaber, E.~Elyan, and C.~Jayne,
``Imitation Learning: A Survey of Learning Methods,''
\emph{ACM Computing Surveys}, vol.~50, no.~2, art.~21, 2017.

\bibitem{vpt2022}
B.~Baker \emph{et al.},
``Video PreTraining (VPT): Learning to Act by Watching Unlabeled Online
Videos,''
in \emph{Advances in Neural Information Processing Systems (NeurIPS)},
vol.~35, 2022.

% -------------------------------------------------------------------------
% 5. Deep Reinforcement Learning -- fundamentos
% -------------------------------------------------------------------------

\bibitem{suttonbarto}
R.~S. Sutton and A.~G. Barto,
\emph{Reinforcement Learning: An Introduction}, 2nd ed.
Cambridge, MA: MIT Press, 2018.

\bibitem{dqnnature}
V.~Mnih \emph{et al.},
``Human-level control through deep reinforcement learning,''
\emph{Nature}, vol.~518, pp.~529--533, 2015.

\bibitem{atari2013}
V.~Mnih, K.~Kavukcuoglu, D.~Silver, A.~Graves, I.~Antonoglou, D.~Wierstra,
and M.~Riedmiller,
``Playing Atari with Deep Reinforcement Learning,''
\emph{arXiv preprint arXiv:1312.5602}, 2013.

\bibitem{doubledqn}
H.~van Hasselt, A.~Guez, and D.~Silver,
``Deep Reinforcement Learning with Double Q-Learning,''
in \emph{Proc.\ 30th AAAI Conf.\ on Artificial Intelligence}, 2016,
pp.~2094--2100.

\bibitem{dueling}
Z.~Wang, T.~Schaul, M.~Hessel, H.~van Hasselt, M.~Lanctot, and N.~de~Freitas,
``Dueling Network Architectures for Deep Reinforcement Learning,''
in \emph{Proc.\ 33rd Int.\ Conf.\ on Machine Learning (ICML)}, 2016.

\bibitem{per}
T.~Schaul, J.~Quan, I.~Antonoglou, and D.~Silver,
``Prioritized Experience Replay,''
in \emph{Proc.\ Int.\ Conf.\ on Learning Representations (ICLR)}, 2016.

\bibitem{rainbow}
M.~Hessel \emph{et al.},
``Rainbow: Combining Improvements in Deep Reinforcement Learning,''
in \emph{Proc.\ 32nd AAAI Conf.\ on Artificial Intelligence}, 2018.

% -------------------------------------------------------------------------
% 6. Policy gradient: PPO, SAC y variantes
% -------------------------------------------------------------------------

\bibitem{ppo}
J.~Schulman, F.~Wolski, P.~Dhariwal, A.~Radford, and O.~Klimov,
``Proximal Policy Optimization Algorithms,''
\emph{arXiv preprint arXiv:1707.06347}, 2017.

\bibitem{sac2018}
T.~Haarnoja, A.~Zhou, P.~Abbeel, and S.~Levine,
``Soft Actor-Critic: Off-Policy Maximum Entropy Deep Reinforcement Learning
with a Stochastic Actor,''
in \emph{Proc.\ 35th Int.\ Conf.\ on Machine Learning (ICML)}, 2018.

\bibitem{sacapp}
T.~Haarnoja \emph{et al.},
``Soft Actor-Critic Algorithms and Applications,''
\emph{arXiv preprint arXiv:1812.05905}, 2018.

\bibitem{sacdiscrete}
P.~Christodoulou,
``Soft Actor-Critic for Discrete Action Settings,''
\emph{arXiv preprint arXiv:1910.07207}, 2019.

\bibitem{a3c}
V.~Mnih \emph{et al.},
``Asynchronous Methods for Deep Reinforcement Learning,''
in \emph{Proc.\ 33rd Int.\ Conf.\ on Machine Learning (ICML)}, 2016.

\bibitem{andrychowicz2021}
M.~Andrychowicz \emph{et al.},
``What Matters for On-Policy Deep Actor-Critic Methods? A Large-Scale
Study,''
in \emph{Proc.\ Int.\ Conf.\ on Learning Representations (ICLR)}, 2021.

% -------------------------------------------------------------------------
% 7. RL + demostraciones (Hybrid, BC + RL)
% -------------------------------------------------------------------------

\bibitem{dqfd}
T.~Hester \emph{et al.},
``Deep Q-Learning from Demonstrations,''
in \emph{Proc.\ 32nd AAAI Conf.\ on Artificial Intelligence}, 2018.

\bibitem{ddpgfd}
M.~Vecerik \emph{et al.},
``Leveraging Demonstrations for Deep Reinforcement Learning on Robotics
Problems with Sparse Rewards,''
\emph{arXiv preprint arXiv:1707.08817}, 2017.

\bibitem{nair2018}
A.~Nair, B.~McGrew, M.~Andrychowicz, W.~Zaremba, and P.~Abbeel,
``Overcoming Exploration in Reinforcement Learning with Demonstrations,''
in \emph{Proc.\ IEEE Int.\ Conf.\ on Robotics and Automation (ICRA)}, 2018.

\bibitem{dapg}
A.~Rajeswaran \emph{et al.},
``Learning Complex Dexterous Manipulation with Deep Reinforcement Learning
and Demonstrations,''
in \emph{Robotics: Science and Systems (RSS)}, 2018.

\bibitem{sqil}
S.~Reddy, A.~D. Dragan, and S.~Levine,
``SQIL: Imitation Learning via Reinforcement Learning with Sparse Rewards,''
in \emph{Proc.\ Int.\ Conf.\ on Learning Representations (ICLR)}, 2020.

\bibitem{iqlearn}
D.~Garg, S.~Chakraborty, C.~Cundy, J.~Song, and S.~Ermon,
``IQ-Learn: Inverse soft-Q Learning for Imitation,''
in \emph{Advances in Neural Information Processing Systems (NeurIPS)},
vol.~34, 2021.

% -------------------------------------------------------------------------
% 8. RL visual, CNN y normalizacion
% -------------------------------------------------------------------------

\bibitem{resnet}
K.~He, X.~Zhang, S.~Ren, and J.~Sun,
``Deep Residual Learning for Image Recognition,''
in \emph{Proc.\ IEEE Conf.\ on Computer Vision and Pattern Recognition
(CVPR)}, 2016, pp.~770--778.

\bibitem{groupnorm}
Y.~Wu and K.~He,
``Group Normalization,''
in \emph{Proc.\ European Conf.\ on Computer Vision (ECCV)}, 2018,
pp.~3--19.

\bibitem{impala}
L.~Espeholt \emph{et al.},
``IMPALA: Scalable Distributed Deep-RL with Importance Weighted
Actor-Learner Architectures,''
in \emph{Proc.\ 35th Int.\ Conf.\ on Machine Learning (ICML)}, 2018.

\bibitem{drqv2}
D.~Yarats, R.~Fergus, A.~Lazaric, and L.~Pinto,
``Mastering Visual Continuous Control: Improved Data-Augmented
Reinforcement Learning,''
\emph{arXiv preprint arXiv:2107.09645}, 2021.

% -------------------------------------------------------------------------
% 9. World models y RL jerarquico (lineas futuras)
% -------------------------------------------------------------------------

\bibitem{worldmodels}
D.~Ha and J.~Schmidhuber,
``World Models,''
in \emph{Advances in Neural Information Processing Systems (NeurIPS)},
2018.

\bibitem{dreamerv3}
D.~Hafner, J.~Pasukonis, J.~Ba, and T.~Lillicrap,
``Mastering Diverse Domains through World Models,''
\emph{arXiv preprint arXiv:2301.04104}, 2023.

\bibitem{muzero}
J.~Schrittwieser \emph{et al.},
``Mastering Atari, Go, Chess and Shogi by Planning with a Learned Model,''
\emph{Nature}, vol.~588, pp.~604--609, 2020.

\bibitem{options}
R.~S. Sutton, D.~Precup, and S.~Singh,
``Between MDPs and semi-MDPs: A Framework for Temporal Abstraction in
Reinforcement Learning,''
\emph{Artificial Intelligence}, vol.~112, no.~1--2, pp.~181--211, 1999.

% -------------------------------------------------------------------------
% 10. Multi-agente
% -------------------------------------------------------------------------

\bibitem{coma}
J.~Foerster, G.~Farquhar, T.~Afouras, N.~Nardelli, and S.~Whiteson,
``Counterfactual Multi-Agent Policy Gradients,''
in \emph{Proc.\ 32nd AAAI Conf.\ on Artificial Intelligence}, 2018.

\bibitem{qmix}
T.~Rashid, M.~Samvelyan, C.~Schroeder de~Witt, G.~Farquhar, J.~Foerster,
and S.~Whiteson,
``QMIX: Monotonic Value Function Factorisation for Deep Multi-Agent
Reinforcement Learning,''
in \emph{Proc.\ 35th Int.\ Conf.\ on Machine Learning (ICML)}, 2018.

\bibitem{maddpg}
R.~Lowe, Y.~Wu, A.~Tamar, J.~Harb, P.~Abbeel, and I.~Mordatch,
``Multi-Agent Actor-Critic for Mixed Cooperative-Competitive
Environments,''
in \emph{Advances in Neural Information Processing Systems (NeurIPS)},
vol.~30, 2017.

% -------------------------------------------------------------------------
% 11. Vision por computador (YOLO, ResNet -- si justificas pipeline visual)
% -------------------------------------------------------------------------

\bibitem{yolo2016}
J.~Redmon, S.~Divvala, R.~Girshick, and A.~Farhadi,
``You Only Look Once: Unified, Real-Time Object Detection,''
in \emph{Proc.\ IEEE Conf.\ on Computer Vision and Pattern Recognition
(CVPR)}, 2016, pp.~779--788.

\bibitem{ultralytics}
G.~Jocher \emph{et al.},
\emph{Ultralytics YOLO} [Software], 2024.
[Online]. Available: \url{https://github.com/ultralytics/ultralytics}

\bibitem{alexnet}
A.~Krizhevsky, I.~Sutskever, and G.~E. Hinton,
``ImageNet Classification with Deep Convolutional Neural Networks,''
in \emph{Advances in Neural Information Processing Systems (NeurIPS)},
vol.~25, 2012, pp.~1097--1105.

\bibitem{convlstm}
X.~Shi, Z.~Chen, H.~Wang, D.-Y. Yeung, W.-K. Wong, and W.-C. Woo,
``Convolutional LSTM Network: A Machine Learning Approach for Precipitation
Nowcasting,''
in \emph{Advances in Neural Information Processing Systems (NeurIPS)},
vol.~28, 2015.

% -------------------------------------------------------------------------
% 12. Surveys y panoramicas (estado del arte)
% -------------------------------------------------------------------------

\bibitem{arulkumaran2017}
K.~Arulkumaran, M.~P. Deisenroth, M.~Brundage, and A.~A. Bharath,
``Deep Reinforcement Learning: A Brief Survey,''
\emph{IEEE Signal Processing Magazine}, vol.~34, no.~6, pp.~26--38, 2017.

\bibitem{li2018}
Y.~Li,
``Deep Reinforcement Learning: An Overview,''
\emph{arXiv preprint arXiv:1701.07274}, 2018.

\bibitem{agentdivisions}
B.~Liu \emph{et al.},
``A Survey of Agent Divisions and Cooperative Frameworks in Multi-Agent
Reinforcement Learning,''
\emph{arXiv preprint arXiv:2111.08857}, 2021.

\end{thebibliography}

\end{document}
   % dentro del documento, antes de \end{document}
%   (c) Copiar solo el bloque \begin{thebibliography}...\end{thebibliography}
%       al final de tu memoria.tex
%
% Estilo: numerico tipo IEEE. Los \cite{key} usan las claves indicadas.
% Entradas marcadas con (*) son las imprescindibles.
% =========================================================================

\documentclass[11pt,a4paper]{article}
\usepackage[utf8]{inputenc}
\usepackage[T1]{fontenc}
\usepackage[spanish]{babel}
\usepackage{geometry}
\usepackage{hyperref}
\geometry{margin=2.5cm}

\title{Bibliografia --- TFG Agentes RL en Minecraft}
\author{Pablo Chantada Saborido}
\date{}

\begin{document}
\maketitle

\begin{thebibliography}{99}

% -------------------------------------------------------------------------
% 1. Plataformas y entornos Minecraft
% -------------------------------------------------------------------------

\bibitem{malmo2016}
M.~Johnson, K.~Hofmann, T.~Hutton, and D.~Bignell,
``The Malmo Platform for Artificial Intelligence Experimentation,''
in \emph{Proc.\ 25th Int.\ Joint Conf.\ on Artificial Intelligence (IJCAI)},
2016, pp.~4246--4247.

\bibitem{minerl2019}
W.~H. Guss, B.~Houghton, N.~Topin, P.~Wang, C.~Codel, M.~Veloso, and
R.~Salakhutdinov,
``MineRL: A Large-Scale Dataset of Minecraft Demonstrations,''
in \emph{Proc.\ 28th Int.\ Joint Conf.\ on Artificial Intelligence (IJCAI)},
2019, pp.~2442--2448.

\bibitem{minerl2020}
W.~H. Guss \emph{et al.},
``Towards Robust and Domain Agnostic Reinforcement Learning Competitions:
MineRL 2020,''
in \emph{Proc.\ NeurIPS 2020 Competition and Demonstration Track},
PMLR, vol.~133, 2021, pp.~233--252.

\bibitem{minerldiamond2021}
A.~Kanervisto \emph{et al.},
``MineRL Diamond 2021 Competition: Overview, Results, and Lessons Learned,''
in \emph{NeurIPS 2021 Competitions and Demonstrations Track},
2022.

\bibitem{minedojo2022}
L.~Fan \emph{et al.},
``MineDojo: Building Open-Ended Embodied Agents with Internet-Scale
Knowledge,''
in \emph{Advances in Neural Information Processing Systems (NeurIPS)},
vol.~35, 2022.

\bibitem{mineflayer}
PrismarineJS,
\emph{Mineflayer: Create Minecraft bots with a high-level JavaScript API}
[Software], 2024. [Online]. Available:
\url{https://github.com/PrismarineJS/mineflayer}

\bibitem{papermc}
PaperMC,
\emph{Paper: High performance Minecraft server} [Software], 2024.
[Online]. Available: \url{https://papermc.io/}

% -------------------------------------------------------------------------
% 2. Agentes con LLMs en Minecraft
% -------------------------------------------------------------------------

\bibitem{voyager2023}
G.~Wang, Y.~Xie, Y.~Jiang, A.~Mandlekar, C.~Xiao, Y.~Zhu, L.~Fan, and
A.~Anandkumar,
``Voyager: An Open-Ended Embodied Agent with Large Language Models,''
\emph{arXiv preprint arXiv:2305.16291}, 2023.

\bibitem{ghost2023}
X.~Zhu \emph{et al.},
``Ghost in the Minecraft: Generally Capable Agents for Open-World Environments
via Large Language Models with Text-based Knowledge and Memory,''
\emph{arXiv preprint arXiv:2305.17144}, 2023.

\bibitem{jarvis1}
Z.~Wang \emph{et al.},
``JARVIS-1: Open-World Multi-task Agents with Memory-Augmented Multimodal
Language Models,''
\emph{arXiv preprint arXiv:2311.05997}, 2023.

\bibitem{steve1}
S.~Lifshitz, K.~Paster, H.~Chan, J.~Ba, and S.~McIlraith,
``STEVE-1: A Generative Model for Text-to-Behavior in Minecraft,''
in \emph{Advances in Neural Information Processing Systems (NeurIPS)},
2023.

\bibitem{plan4mc}
H.~Yuan, C.~Zhang, H.~Wang, F.~Xie, P.~Cai, H.~Dong, and Z.~Lu,
``Plan4MC: Skill Reinforcement Learning and Planning for Open-World Minecraft
Tasks,''
\emph{arXiv preprint arXiv:2303.16563}, 2023.

\bibitem{mindcraft}
K.~Nottingham \emph{et al.},
\emph{Mindcraft: LLM-powered Minecraft agents} [Software], 2024.
[Online]. Available: \url{https://github.com/kolbytn/mindcraft}

% -------------------------------------------------------------------------
% 3. Planificacion clasica (HTN)
% -------------------------------------------------------------------------

\bibitem{erol1994}
K.~Erol, J.~Hendler, and D.~S. Nau,
``HTN Planning: Complexity and Expressivity,''
in \emph{Proc.\ 12th National Conf.\ on Artificial Intelligence (AAAI)},
1994, pp.~1123--1128.

\bibitem{ghallab2004}
M.~Ghallab, D.~Nau, and P.~Traverso,
\emph{Automated Planning: Theory and Practice}.
San Francisco, CA: Morgan Kaufmann, 2004.

\bibitem{shop2}
D.~S. Nau, T.-C. Au, O.~Ilghami, U.~Kuter, J.~W. Murdock, D.~Wu, and F.~Yaman,
``SHOP2: An HTN Planning System,''
\emph{Journal of Artificial Intelligence Research}, vol.~20, pp.~379--404,
2003.

\bibitem{georgievski2015}
I.~Georgievski and M.~Aiello,
``HTN planning: Overview, comparison, and beyond,''
\emph{Artificial Intelligence}, vol.~222, pp.~124--156, 2015.

% -------------------------------------------------------------------------
% 4. Imitation Learning y Behavior Cloning
% -------------------------------------------------------------------------

\bibitem{pomerleau1991}
D.~A. Pomerleau,
``Efficient Training of Artificial Neural Networks for Autonomous
Navigation,''
\emph{Neural Computation}, vol.~3, no.~1, pp.~88--97, 1991.

\bibitem{dagger}
S.~Ross, G.~Gordon, and D.~Bagnell,
``A Reduction of Imitation Learning and Structured Prediction to No-Regret
Online Learning,''
in \emph{Proc.\ 14th Int.\ Conf.\ on Artificial Intelligence and Statistics
(AISTATS)}, 2011, pp.~627--635.

\bibitem{gail}
J.~Ho and S.~Ermon,
``Generative Adversarial Imitation Learning,''
in \emph{Advances in Neural Information Processing Systems (NeurIPS)},
vol.~29, 2016.

\bibitem{hussein2017}
A.~Hussein, M.~M. Gaber, E.~Elyan, and C.~Jayne,
``Imitation Learning: A Survey of Learning Methods,''
\emph{ACM Computing Surveys}, vol.~50, no.~2, art.~21, 2017.

\bibitem{vpt2022}
B.~Baker \emph{et al.},
``Video PreTraining (VPT): Learning to Act by Watching Unlabeled Online
Videos,''
in \emph{Advances in Neural Information Processing Systems (NeurIPS)},
vol.~35, 2022.

% -------------------------------------------------------------------------
% 5. Deep Reinforcement Learning -- fundamentos
% -------------------------------------------------------------------------

\bibitem{suttonbarto}
R.~S. Sutton and A.~G. Barto,
\emph{Reinforcement Learning: An Introduction}, 2nd ed.
Cambridge, MA: MIT Press, 2018.

\bibitem{dqnnature}
V.~Mnih \emph{et al.},
``Human-level control through deep reinforcement learning,''
\emph{Nature}, vol.~518, pp.~529--533, 2015.

\bibitem{atari2013}
V.~Mnih, K.~Kavukcuoglu, D.~Silver, A.~Graves, I.~Antonoglou, D.~Wierstra,
and M.~Riedmiller,
``Playing Atari with Deep Reinforcement Learning,''
\emph{arXiv preprint arXiv:1312.5602}, 2013.

\bibitem{doubledqn}
H.~van Hasselt, A.~Guez, and D.~Silver,
``Deep Reinforcement Learning with Double Q-Learning,''
in \emph{Proc.\ 30th AAAI Conf.\ on Artificial Intelligence}, 2016,
pp.~2094--2100.

\bibitem{dueling}
Z.~Wang, T.~Schaul, M.~Hessel, H.~van Hasselt, M.~Lanctot, and N.~de~Freitas,
``Dueling Network Architectures for Deep Reinforcement Learning,''
in \emph{Proc.\ 33rd Int.\ Conf.\ on Machine Learning (ICML)}, 2016.

\bibitem{per}
T.~Schaul, J.~Quan, I.~Antonoglou, and D.~Silver,
``Prioritized Experience Replay,''
in \emph{Proc.\ Int.\ Conf.\ on Learning Representations (ICLR)}, 2016.

\bibitem{rainbow}
M.~Hessel \emph{et al.},
``Rainbow: Combining Improvements in Deep Reinforcement Learning,''
in \emph{Proc.\ 32nd AAAI Conf.\ on Artificial Intelligence}, 2018.

% -------------------------------------------------------------------------
% 6. Policy gradient: PPO, SAC y variantes
% -------------------------------------------------------------------------

\bibitem{ppo}
J.~Schulman, F.~Wolski, P.~Dhariwal, A.~Radford, and O.~Klimov,
``Proximal Policy Optimization Algorithms,''
\emph{arXiv preprint arXiv:1707.06347}, 2017.

\bibitem{sac2018}
T.~Haarnoja, A.~Zhou, P.~Abbeel, and S.~Levine,
``Soft Actor-Critic: Off-Policy Maximum Entropy Deep Reinforcement Learning
with a Stochastic Actor,''
in \emph{Proc.\ 35th Int.\ Conf.\ on Machine Learning (ICML)}, 2018.

\bibitem{sacapp}
T.~Haarnoja \emph{et al.},
``Soft Actor-Critic Algorithms and Applications,''
\emph{arXiv preprint arXiv:1812.05905}, 2018.

\bibitem{sacdiscrete}
P.~Christodoulou,
``Soft Actor-Critic for Discrete Action Settings,''
\emph{arXiv preprint arXiv:1910.07207}, 2019.

\bibitem{a3c}
V.~Mnih \emph{et al.},
``Asynchronous Methods for Deep Reinforcement Learning,''
in \emph{Proc.\ 33rd Int.\ Conf.\ on Machine Learning (ICML)}, 2016.

\bibitem{andrychowicz2021}
M.~Andrychowicz \emph{et al.},
``What Matters for On-Policy Deep Actor-Critic Methods? A Large-Scale
Study,''
in \emph{Proc.\ Int.\ Conf.\ on Learning Representations (ICLR)}, 2021.

% -------------------------------------------------------------------------
% 7. RL + demostraciones (Hybrid, BC + RL)
% -------------------------------------------------------------------------

\bibitem{dqfd}
T.~Hester \emph{et al.},
``Deep Q-Learning from Demonstrations,''
in \emph{Proc.\ 32nd AAAI Conf.\ on Artificial Intelligence}, 2018.

\bibitem{ddpgfd}
M.~Vecerik \emph{et al.},
``Leveraging Demonstrations for Deep Reinforcement Learning on Robotics
Problems with Sparse Rewards,''
\emph{arXiv preprint arXiv:1707.08817}, 2017.

\bibitem{nair2018}
A.~Nair, B.~McGrew, M.~Andrychowicz, W.~Zaremba, and P.~Abbeel,
``Overcoming Exploration in Reinforcement Learning with Demonstrations,''
in \emph{Proc.\ IEEE Int.\ Conf.\ on Robotics and Automation (ICRA)}, 2018.

\bibitem{dapg}
A.~Rajeswaran \emph{et al.},
``Learning Complex Dexterous Manipulation with Deep Reinforcement Learning
and Demonstrations,''
in \emph{Robotics: Science and Systems (RSS)}, 2018.

\bibitem{sqil}
S.~Reddy, A.~D. Dragan, and S.~Levine,
``SQIL: Imitation Learning via Reinforcement Learning with Sparse Rewards,''
in \emph{Proc.\ Int.\ Conf.\ on Learning Representations (ICLR)}, 2020.

\bibitem{iqlearn}
D.~Garg, S.~Chakraborty, C.~Cundy, J.~Song, and S.~Ermon,
``IQ-Learn: Inverse soft-Q Learning for Imitation,''
in \emph{Advances in Neural Information Processing Systems (NeurIPS)},
vol.~34, 2021.

% -------------------------------------------------------------------------
% 8. RL visual, CNN y normalizacion
% -------------------------------------------------------------------------

\bibitem{resnet}
K.~He, X.~Zhang, S.~Ren, and J.~Sun,
``Deep Residual Learning for Image Recognition,''
in \emph{Proc.\ IEEE Conf.\ on Computer Vision and Pattern Recognition
(CVPR)}, 2016, pp.~770--778.

\bibitem{groupnorm}
Y.~Wu and K.~He,
``Group Normalization,''
in \emph{Proc.\ European Conf.\ on Computer Vision (ECCV)}, 2018,
pp.~3--19.

\bibitem{impala}
L.~Espeholt \emph{et al.},
``IMPALA: Scalable Distributed Deep-RL with Importance Weighted
Actor-Learner Architectures,''
in \emph{Proc.\ 35th Int.\ Conf.\ on Machine Learning (ICML)}, 2018.

\bibitem{drqv2}
D.~Yarats, R.~Fergus, A.~Lazaric, and L.~Pinto,
``Mastering Visual Continuous Control: Improved Data-Augmented
Reinforcement Learning,''
\emph{arXiv preprint arXiv:2107.09645}, 2021.

% -------------------------------------------------------------------------
% 9. World models y RL jerarquico (lineas futuras)
% -------------------------------------------------------------------------

\bibitem{worldmodels}
D.~Ha and J.~Schmidhuber,
``World Models,''
in \emph{Advances in Neural Information Processing Systems (NeurIPS)},
2018.

\bibitem{dreamerv3}
D.~Hafner, J.~Pasukonis, J.~Ba, and T.~Lillicrap,
``Mastering Diverse Domains through World Models,''
\emph{arXiv preprint arXiv:2301.04104}, 2023.

\bibitem{muzero}
J.~Schrittwieser \emph{et al.},
``Mastering Atari, Go, Chess and Shogi by Planning with a Learned Model,''
\emph{Nature}, vol.~588, pp.~604--609, 2020.

\bibitem{options}
R.~S. Sutton, D.~Precup, and S.~Singh,
``Between MDPs and semi-MDPs: A Framework for Temporal Abstraction in
Reinforcement Learning,''
\emph{Artificial Intelligence}, vol.~112, no.~1--2, pp.~181--211, 1999.

% -------------------------------------------------------------------------
% 10. Multi-agente
% -------------------------------------------------------------------------

\bibitem{coma}
J.~Foerster, G.~Farquhar, T.~Afouras, N.~Nardelli, and S.~Whiteson,
``Counterfactual Multi-Agent Policy Gradients,''
in \emph{Proc.\ 32nd AAAI Conf.\ on Artificial Intelligence}, 2018.

\bibitem{qmix}
T.~Rashid, M.~Samvelyan, C.~Schroeder de~Witt, G.~Farquhar, J.~Foerster,
and S.~Whiteson,
``QMIX: Monotonic Value Function Factorisation for Deep Multi-Agent
Reinforcement Learning,''
in \emph{Proc.\ 35th Int.\ Conf.\ on Machine Learning (ICML)}, 2018.

\bibitem{maddpg}
R.~Lowe, Y.~Wu, A.~Tamar, J.~Harb, P.~Abbeel, and I.~Mordatch,
``Multi-Agent Actor-Critic for Mixed Cooperative-Competitive
Environments,''
in \emph{Advances in Neural Information Processing Systems (NeurIPS)},
vol.~30, 2017.

% -------------------------------------------------------------------------
% 11. Vision por computador (YOLO, ResNet -- si justificas pipeline visual)
% -------------------------------------------------------------------------

\bibitem{yolo2016}
J.~Redmon, S.~Divvala, R.~Girshick, and A.~Farhadi,
``You Only Look Once: Unified, Real-Time Object Detection,''
in \emph{Proc.\ IEEE Conf.\ on Computer Vision and Pattern Recognition
(CVPR)}, 2016, pp.~779--788.

\bibitem{ultralytics}
G.~Jocher \emph{et al.},
\emph{Ultralytics YOLO} [Software], 2024.
[Online]. Available: \url{https://github.com/ultralytics/ultralytics}

\bibitem{alexnet}
A.~Krizhevsky, I.~Sutskever, and G.~E. Hinton,
``ImageNet Classification with Deep Convolutional Neural Networks,''
in \emph{Advances in Neural Information Processing Systems (NeurIPS)},
vol.~25, 2012, pp.~1097--1105.

\bibitem{convlstm}
X.~Shi, Z.~Chen, H.~Wang, D.-Y. Yeung, W.-K. Wong, and W.-C. Woo,
``Convolutional LSTM Network: A Machine Learning Approach for Precipitation
Nowcasting,''
in \emph{Advances in Neural Information Processing Systems (NeurIPS)},
vol.~28, 2015.

% -------------------------------------------------------------------------
% 12. Surveys y panoramicas (estado del arte)
% -------------------------------------------------------------------------

\bibitem{arulkumaran2017}
K.~Arulkumaran, M.~P. Deisenroth, M.~Brundage, and A.~A. Bharath,
``Deep Reinforcement Learning: A Brief Survey,''
\emph{IEEE Signal Processing Magazine}, vol.~34, no.~6, pp.~26--38, 2017.

\bibitem{li2018}
Y.~Li,
``Deep Reinforcement Learning: An Overview,''
\emph{arXiv preprint arXiv:1701.07274}, 2018.

\bibitem{agentdivisions}
B.~Liu \emph{et al.},
``A Survey of Agent Divisions and Cooperative Frameworks in Multi-Agent
Reinforcement Learning,''
\emph{arXiv preprint arXiv:2111.08857}, 2021.

\end{thebibliography}

\end{document}
   % dentro del documento, antes de \end{document}
%   (c) Copiar solo el bloque \begin{thebibliography}...\end{thebibliography}
%       al final de tu memoria.tex
%
% Estilo: numerico tipo IEEE. Los \cite{key} usan las claves indicadas.
% Entradas marcadas con (*) son las imprescindibles.
% =========================================================================

\documentclass[11pt,a4paper]{article}
\usepackage[utf8]{inputenc}
\usepackage[T1]{fontenc}
\usepackage[spanish]{babel}
\usepackage{geometry}
\usepackage{hyperref}
\geometry{margin=2.5cm}

\title{Bibliografia --- TFG Agentes RL en Minecraft}
\author{Pablo Chantada Saborido}
\date{}

\begin{document}
\maketitle

\begin{thebibliography}{99}

% -------------------------------------------------------------------------
% 1. Plataformas y entornos Minecraft
% -------------------------------------------------------------------------

\bibitem{malmo2016}
M.~Johnson, K.~Hofmann, T.~Hutton, and D.~Bignell,
``The Malmo Platform for Artificial Intelligence Experimentation,''
in \emph{Proc.\ 25th Int.\ Joint Conf.\ on Artificial Intelligence (IJCAI)},
2016, pp.~4246--4247.

\bibitem{minerl2019}
W.~H. Guss, B.~Houghton, N.~Topin, P.~Wang, C.~Codel, M.~Veloso, and
R.~Salakhutdinov,
``MineRL: A Large-Scale Dataset of Minecraft Demonstrations,''
in \emph{Proc.\ 28th Int.\ Joint Conf.\ on Artificial Intelligence (IJCAI)},
2019, pp.~2442--2448.

\bibitem{minerl2020}
W.~H. Guss \emph{et al.},
``Towards Robust and Domain Agnostic Reinforcement Learning Competitions:
MineRL 2020,''
in \emph{Proc.\ NeurIPS 2020 Competition and Demonstration Track},
PMLR, vol.~133, 2021, pp.~233--252.

\bibitem{minerldiamond2021}
A.~Kanervisto \emph{et al.},
``MineRL Diamond 2021 Competition: Overview, Results, and Lessons Learned,''
in \emph{NeurIPS 2021 Competitions and Demonstrations Track},
2022.

\bibitem{minedojo2022}
L.~Fan \emph{et al.},
``MineDojo: Building Open-Ended Embodied Agents with Internet-Scale
Knowledge,''
in \emph{Advances in Neural Information Processing Systems (NeurIPS)},
vol.~35, 2022.

\bibitem{mineflayer}
PrismarineJS,
\emph{Mineflayer: Create Minecraft bots with a high-level JavaScript API}
[Software], 2024. [Online]. Available:
\url{https://github.com/PrismarineJS/mineflayer}

\bibitem{papermc}
PaperMC,
\emph{Paper: High performance Minecraft server} [Software], 2024.
[Online]. Available: \url{https://papermc.io/}

% -------------------------------------------------------------------------
% 2. Agentes con LLMs en Minecraft
% -------------------------------------------------------------------------

\bibitem{voyager2023}
G.~Wang, Y.~Xie, Y.~Jiang, A.~Mandlekar, C.~Xiao, Y.~Zhu, L.~Fan, and
A.~Anandkumar,
``Voyager: An Open-Ended Embodied Agent with Large Language Models,''
\emph{arXiv preprint arXiv:2305.16291}, 2023.

\bibitem{ghost2023}
X.~Zhu \emph{et al.},
``Ghost in the Minecraft: Generally Capable Agents for Open-World Environments
via Large Language Models with Text-based Knowledge and Memory,''
\emph{arXiv preprint arXiv:2305.17144}, 2023.

\bibitem{jarvis1}
Z.~Wang \emph{et al.},
``JARVIS-1: Open-World Multi-task Agents with Memory-Augmented Multimodal
Language Models,''
\emph{arXiv preprint arXiv:2311.05997}, 2023.

\bibitem{steve1}
S.~Lifshitz, K.~Paster, H.~Chan, J.~Ba, and S.~McIlraith,
``STEVE-1: A Generative Model for Text-to-Behavior in Minecraft,''
in \emph{Advances in Neural Information Processing Systems (NeurIPS)},
2023.

\bibitem{plan4mc}
H.~Yuan, C.~Zhang, H.~Wang, F.~Xie, P.~Cai, H.~Dong, and Z.~Lu,
``Plan4MC: Skill Reinforcement Learning and Planning for Open-World Minecraft
Tasks,''
\emph{arXiv preprint arXiv:2303.16563}, 2023.

\bibitem{mindcraft}
K.~Nottingham \emph{et al.},
\emph{Mindcraft: LLM-powered Minecraft agents} [Software], 2024.
[Online]. Available: \url{https://github.com/kolbytn/mindcraft}

% -------------------------------------------------------------------------
% 3. Planificacion clasica (HTN)
% -------------------------------------------------------------------------

\bibitem{erol1994}
K.~Erol, J.~Hendler, and D.~S. Nau,
``HTN Planning: Complexity and Expressivity,''
in \emph{Proc.\ 12th National Conf.\ on Artificial Intelligence (AAAI)},
1994, pp.~1123--1128.

\bibitem{ghallab2004}
M.~Ghallab, D.~Nau, and P.~Traverso,
\emph{Automated Planning: Theory and Practice}.
San Francisco, CA: Morgan Kaufmann, 2004.

\bibitem{shop2}
D.~S. Nau, T.-C. Au, O.~Ilghami, U.~Kuter, J.~W. Murdock, D.~Wu, and F.~Yaman,
``SHOP2: An HTN Planning System,''
\emph{Journal of Artificial Intelligence Research}, vol.~20, pp.~379--404,
2003.

\bibitem{georgievski2015}
I.~Georgievski and M.~Aiello,
``HTN planning: Overview, comparison, and beyond,''
\emph{Artificial Intelligence}, vol.~222, pp.~124--156, 2015.

% -------------------------------------------------------------------------
% 4. Imitation Learning y Behavior Cloning
% -------------------------------------------------------------------------

\bibitem{pomerleau1991}
D.~A. Pomerleau,
``Efficient Training of Artificial Neural Networks for Autonomous
Navigation,''
\emph{Neural Computation}, vol.~3, no.~1, pp.~88--97, 1991.

\bibitem{dagger}
S.~Ross, G.~Gordon, and D.~Bagnell,
``A Reduction of Imitation Learning and Structured Prediction to No-Regret
Online Learning,''
in \emph{Proc.\ 14th Int.\ Conf.\ on Artificial Intelligence and Statistics
(AISTATS)}, 2011, pp.~627--635.

\bibitem{gail}
J.~Ho and S.~Ermon,
``Generative Adversarial Imitation Learning,''
in \emph{Advances in Neural Information Processing Systems (NeurIPS)},
vol.~29, 2016.

\bibitem{hussein2017}
A.~Hussein, M.~M. Gaber, E.~Elyan, and C.~Jayne,
``Imitation Learning: A Survey of Learning Methods,''
\emph{ACM Computing Surveys}, vol.~50, no.~2, art.~21, 2017.

\bibitem{vpt2022}
B.~Baker \emph{et al.},
``Video PreTraining (VPT): Learning to Act by Watching Unlabeled Online
Videos,''
in \emph{Advances in Neural Information Processing Systems (NeurIPS)},
vol.~35, 2022.

% -------------------------------------------------------------------------
% 5. Deep Reinforcement Learning -- fundamentos
% -------------------------------------------------------------------------

\bibitem{suttonbarto}
R.~S. Sutton and A.~G. Barto,
\emph{Reinforcement Learning: An Introduction}, 2nd ed.
Cambridge, MA: MIT Press, 2018.

\bibitem{dqnnature}
V.~Mnih \emph{et al.},
``Human-level control through deep reinforcement learning,''
\emph{Nature}, vol.~518, pp.~529--533, 2015.

\bibitem{atari2013}
V.~Mnih, K.~Kavukcuoglu, D.~Silver, A.~Graves, I.~Antonoglou, D.~Wierstra,
and M.~Riedmiller,
``Playing Atari with Deep Reinforcement Learning,''
\emph{arXiv preprint arXiv:1312.5602}, 2013.

\bibitem{doubledqn}
H.~van Hasselt, A.~Guez, and D.~Silver,
``Deep Reinforcement Learning with Double Q-Learning,''
in \emph{Proc.\ 30th AAAI Conf.\ on Artificial Intelligence}, 2016,
pp.~2094--2100.

\bibitem{dueling}
Z.~Wang, T.~Schaul, M.~Hessel, H.~van Hasselt, M.~Lanctot, and N.~de~Freitas,
``Dueling Network Architectures for Deep Reinforcement Learning,''
in \emph{Proc.\ 33rd Int.\ Conf.\ on Machine Learning (ICML)}, 2016.

\bibitem{per}
T.~Schaul, J.~Quan, I.~Antonoglou, and D.~Silver,
``Prioritized Experience Replay,''
in \emph{Proc.\ Int.\ Conf.\ on Learning Representations (ICLR)}, 2016.

\bibitem{rainbow}
M.~Hessel \emph{et al.},
``Rainbow: Combining Improvements in Deep Reinforcement Learning,''
in \emph{Proc.\ 32nd AAAI Conf.\ on Artificial Intelligence}, 2018.

% -------------------------------------------------------------------------
% 6. Policy gradient: PPO, SAC y variantes
% -------------------------------------------------------------------------

\bibitem{ppo}
J.~Schulman, F.~Wolski, P.~Dhariwal, A.~Radford, and O.~Klimov,
``Proximal Policy Optimization Algorithms,''
\emph{arXiv preprint arXiv:1707.06347}, 2017.

\bibitem{sac2018}
T.~Haarnoja, A.~Zhou, P.~Abbeel, and S.~Levine,
``Soft Actor-Critic: Off-Policy Maximum Entropy Deep Reinforcement Learning
with a Stochastic Actor,''
in \emph{Proc.\ 35th Int.\ Conf.\ on Machine Learning (ICML)}, 2018.

\bibitem{sacapp}
T.~Haarnoja \emph{et al.},
``Soft Actor-Critic Algorithms and Applications,''
\emph{arXiv preprint arXiv:1812.05905}, 2018.

\bibitem{sacdiscrete}
P.~Christodoulou,
``Soft Actor-Critic for Discrete Action Settings,''
\emph{arXiv preprint arXiv:1910.07207}, 2019.

\bibitem{a3c}
V.~Mnih \emph{et al.},
``Asynchronous Methods for Deep Reinforcement Learning,''
in \emph{Proc.\ 33rd Int.\ Conf.\ on Machine Learning (ICML)}, 2016.

\bibitem{andrychowicz2021}
M.~Andrychowicz \emph{et al.},
``What Matters for On-Policy Deep Actor-Critic Methods? A Large-Scale
Study,''
in \emph{Proc.\ Int.\ Conf.\ on Learning Representations (ICLR)}, 2021.

% -------------------------------------------------------------------------
% 7. RL + demostraciones (Hybrid, BC + RL)
% -------------------------------------------------------------------------

\bibitem{dqfd}
T.~Hester \emph{et al.},
``Deep Q-Learning from Demonstrations,''
in \emph{Proc.\ 32nd AAAI Conf.\ on Artificial Intelligence}, 2018.

\bibitem{ddpgfd}
M.~Vecerik \emph{et al.},
``Leveraging Demonstrations for Deep Reinforcement Learning on Robotics
Problems with Sparse Rewards,''
\emph{arXiv preprint arXiv:1707.08817}, 2017.

\bibitem{nair2018}
A.~Nair, B.~McGrew, M.~Andrychowicz, W.~Zaremba, and P.~Abbeel,
``Overcoming Exploration in Reinforcement Learning with Demonstrations,''
in \emph{Proc.\ IEEE Int.\ Conf.\ on Robotics and Automation (ICRA)}, 2018.

\bibitem{dapg}
A.~Rajeswaran \emph{et al.},
``Learning Complex Dexterous Manipulation with Deep Reinforcement Learning
and Demonstrations,''
in \emph{Robotics: Science and Systems (RSS)}, 2018.

\bibitem{sqil}
S.~Reddy, A.~D. Dragan, and S.~Levine,
``SQIL: Imitation Learning via Reinforcement Learning with Sparse Rewards,''
in \emph{Proc.\ Int.\ Conf.\ on Learning Representations (ICLR)}, 2020.

\bibitem{iqlearn}
D.~Garg, S.~Chakraborty, C.~Cundy, J.~Song, and S.~Ermon,
``IQ-Learn: Inverse soft-Q Learning for Imitation,''
in \emph{Advances in Neural Information Processing Systems (NeurIPS)},
vol.~34, 2021.

% -------------------------------------------------------------------------
% 8. RL visual, CNN y normalizacion
% -------------------------------------------------------------------------

\bibitem{resnet}
K.~He, X.~Zhang, S.~Ren, and J.~Sun,
``Deep Residual Learning for Image Recognition,''
in \emph{Proc.\ IEEE Conf.\ on Computer Vision and Pattern Recognition
(CVPR)}, 2016, pp.~770--778.

\bibitem{groupnorm}
Y.~Wu and K.~He,
``Group Normalization,''
in \emph{Proc.\ European Conf.\ on Computer Vision (ECCV)}, 2018,
pp.~3--19.

\bibitem{impala}
L.~Espeholt \emph{et al.},
``IMPALA: Scalable Distributed Deep-RL with Importance Weighted
Actor-Learner Architectures,''
in \emph{Proc.\ 35th Int.\ Conf.\ on Machine Learning (ICML)}, 2018.

\bibitem{drqv2}
D.~Yarats, R.~Fergus, A.~Lazaric, and L.~Pinto,
``Mastering Visual Continuous Control: Improved Data-Augmented
Reinforcement Learning,''
\emph{arXiv preprint arXiv:2107.09645}, 2021.

% -------------------------------------------------------------------------
% 9. World models y RL jerarquico (lineas futuras)
% -------------------------------------------------------------------------

\bibitem{worldmodels}
D.~Ha and J.~Schmidhuber,
``World Models,''
in \emph{Advances in Neural Information Processing Systems (NeurIPS)},
2018.

\bibitem{dreamerv3}
D.~Hafner, J.~Pasukonis, J.~Ba, and T.~Lillicrap,
``Mastering Diverse Domains through World Models,''
\emph{arXiv preprint arXiv:2301.04104}, 2023.

\bibitem{muzero}
J.~Schrittwieser \emph{et al.},
``Mastering Atari, Go, Chess and Shogi by Planning with a Learned Model,''
\emph{Nature}, vol.~588, pp.~604--609, 2020.

\bibitem{options}
R.~S. Sutton, D.~Precup, and S.~Singh,
``Between MDPs and semi-MDPs: A Framework for Temporal Abstraction in
Reinforcement Learning,''
\emph{Artificial Intelligence}, vol.~112, no.~1--2, pp.~181--211, 1999.

% -------------------------------------------------------------------------
% 10. Multi-agente
% -------------------------------------------------------------------------

\bibitem{coma}
J.~Foerster, G.~Farquhar, T.~Afouras, N.~Nardelli, and S.~Whiteson,
``Counterfactual Multi-Agent Policy Gradients,''
in \emph{Proc.\ 32nd AAAI Conf.\ on Artificial Intelligence}, 2018.

\bibitem{qmix}
T.~Rashid, M.~Samvelyan, C.~Schroeder de~Witt, G.~Farquhar, J.~Foerster,
and S.~Whiteson,
``QMIX: Monotonic Value Function Factorisation for Deep Multi-Agent
Reinforcement Learning,''
in \emph{Proc.\ 35th Int.\ Conf.\ on Machine Learning (ICML)}, 2018.

\bibitem{maddpg}
R.~Lowe, Y.~Wu, A.~Tamar, J.~Harb, P.~Abbeel, and I.~Mordatch,
``Multi-Agent Actor-Critic for Mixed Cooperative-Competitive
Environments,''
in \emph{Advances in Neural Information Processing Systems (NeurIPS)},
vol.~30, 2017.

% -------------------------------------------------------------------------
% 11. Vision por computador (YOLO, ResNet -- si justificas pipeline visual)
% -------------------------------------------------------------------------

\bibitem{yolo2016}
J.~Redmon, S.~Divvala, R.~Girshick, and A.~Farhadi,
``You Only Look Once: Unified, Real-Time Object Detection,''
in \emph{Proc.\ IEEE Conf.\ on Computer Vision and Pattern Recognition
(CVPR)}, 2016, pp.~779--788.

\bibitem{ultralytics}
G.~Jocher \emph{et al.},
\emph{Ultralytics YOLO} [Software], 2024.
[Online]. Available: \url{https://github.com/ultralytics/ultralytics}

\bibitem{alexnet}
A.~Krizhevsky, I.~Sutskever, and G.~E. Hinton,
``ImageNet Classification with Deep Convolutional Neural Networks,''
in \emph{Advances in Neural Information Processing Systems (NeurIPS)},
vol.~25, 2012, pp.~1097--1105.

\bibitem{convlstm}
X.~Shi, Z.~Chen, H.~Wang, D.-Y. Yeung, W.-K. Wong, and W.-C. Woo,
``Convolutional LSTM Network: A Machine Learning Approach for Precipitation
Nowcasting,''
in \emph{Advances in Neural Information Processing Systems (NeurIPS)},
vol.~28, 2015.

% -------------------------------------------------------------------------
% 12. Surveys y panoramicas (estado del arte)
% -------------------------------------------------------------------------

\bibitem{arulkumaran2017}
K.~Arulkumaran, M.~P. Deisenroth, M.~Brundage, and A.~A. Bharath,
``Deep Reinforcement Learning: A Brief Survey,''
\emph{IEEE Signal Processing Magazine}, vol.~34, no.~6, pp.~26--38, 2017.

\bibitem{li2018}
Y.~Li,
``Deep Reinforcement Learning: An Overview,''
\emph{arXiv preprint arXiv:1701.07274}, 2018.

\bibitem{agentdivisions}
B.~Liu \emph{et al.},
``A Survey of Agent Divisions and Cooperative Frameworks in Multi-Agent
Reinforcement Learning,''
\emph{arXiv preprint arXiv:2111.08857}, 2021.

\end{thebibliography}

\end{document}
